\begin{abstract}
\textbf{Background:} Spaceflight exposes biological systems to the synergistic stresses of microgravity, cosmic ionizing radiation, and sealed atmospheric confinement. While the transcriptional responses of model plants like \textit{Arabidopsis thaliana} have been comprehensively cataloged in the NASA Open Science Data Repository (OSDR), the endophytic and epiphytic microbiomes associated with these seedlings have remained virtually unexplored in archival transcriptomic datasets. Because standard eukaryotic RNA sequencing protocols capture both host and microbial polyadenylated or co-extracted transcripts, mining unmapped non-host reads provides a powerful window into plant-microbiome dynamics in extraterrestrial environments.

\textbf{Methods:} Here, we present a systematic meta-transcriptomic extraction and taxonomic profiling pipeline applied across six OSDR spaceflight accessions (OSD-37, OSD-38, OSD-120, OSD-217, OSD-321, and OSD-522) encompassing multiple hardware systems (BRIC-19, BRIC-20, APEX-03-2, BRIC-PDF) and genetic backgrounds (Columbia-0, Wassilewskija-2, Landsberg \textit{erecta}-0, and Cape Verde Islands-0). Following host read subtraction against the \textit{A. thaliana} TAIR10 reference genome, remaining microbial sequences were profiled using MetaPhlAn4, Bowtie2, and Krona interactive visualization. Taxonomic intersections across flight and ground controls were assessed through UpSetR and hierarchical correlation analyses, and annotated using The Microbe Directory.

\textbf{Results:} We observed distinct microbial communities segregated between intrinsic seed-borne endophytes and extrinsic hardware- or human-associated commensals. \textit{Ralstonia pickettii} exhibited pervasive dominance across OSD-37 ecotypes, reaching 100\% relative abundance in Cvi-0 samples. In OSD-38, we uncovered the presence of \textit{Weizmannia ginsengihumi}, a beneficial plant-growth-promoting Firmicute previously unrecorded in orbital seedling habitats. Furthermore, common skin and environmental commensals including \textit{Cutibacterium acnes}, \textit{Malassezia restricta}, and \textit{Escherichia coli} were systematically identified, reflecting microbial exchange between crew-handled hardware and plant growth substrates. Interestingly, comparison of whole seedlings (OSD-37, 38) against dissected root tips (OSD-120) confirmed that endophytic colonization is strongly tissue-dependent, with root tips displaying near-complete depletion of bacterial transcripts.

\textbf{Conclusions:} These findings demonstrate that plant-associated microbiomes persist in closed spaceflight hardware and undergo community shifts that mirror both host developmental state and chamber microbial ecologies. Archival metatranscriptomics provides a cost-effective, non-destructive paradigm for tracking plant-microbe biosecurity, designing beneficial synthetic microbiomes, and ensuring crop resilience in future Bioregenerative Life Support Systems (BLSS) on Lunar and Martian surface missions.
\end{abstract}
